% Kill-Chain Canaries — arXiv submission
% Compile:  pdflatex main  (twice)
% Requires: icml2025.sty (included); figures/killchain_figures/output/*.pdf
% Figures:  cd figures/killchain_figures && python build_all.py
\documentclass{article}

\PassOptionsToPackage{numbers}{natbib} % force numbered citations before icml2025.sty loads natbib
\usepackage{icml2025}          % custom sty — handles geometry, fonts, sections
\usepackage[utf8]{inputenc}
\usepackage{placeins}          % \FloatBarrier — fixes float/section ordering
\usepackage{url}
\usepackage{booktabs}
\usepackage{amsmath}
\usepackage{amssymb}
\usepackage{graphicx}
\usepackage{caption}
\usepackage{subcaption}
\usepackage{xcolor}
\usepackage{listings}
\usepackage{microtype}
\usepackage{multirow}
\usepackage{array}

% Prevent bad mid-word breaks for key proper nouns
\hyphenation{PropagationLogger MemoryStore ToolRegistry MultiAgentOrchestrator
             crypto-graphic sum-mar-iza-tion}

% Code literals.  \path (url package) may break a literal after _ . @ / , =
% so long identifiers never run into the margin; \urldef makes the macros
% safe in captions and tables.  Use them with {} before a following space.
\urldef{\mempoison}\path{memory_poison}
\urldef{\toolpoison}\path{tool_poison}
\urldef{\propag}\path{propagation}
\urldef{\permesc}\path{permission_esc}
\urldef{\multisurf}\path{multi_surface}
\urldef{\crossmodal}\path{cross_modal_relay}
\urldef{\pdfappend}\path{pdf_append}
\urldef{\pdfwhite}\path{pdf_whitefont}
\urldef{\pdfmeta}\path{pdf_metadata}
\urldef{\audioinj}\path{audio_inject}
\urldef{\wfilter}\path{write_filter}
\urldef{\pidetect}\path{pi_detector}
\urldef{\spotlight}\path{spotlighting}
\urldef{\wmem}\path{write_memory}
\urldef{\rmem}\path{read_memory}
\urldef{\sendrep}\path{send_report}
\urldef{\parsepdf}\path{parse_pdf}
\urldef{\getweb}\path{get_webpage}
\urldef{\querydb}\path{query_db}
\urldef{\transcribe}\path{transcribe_audio}
\urldef{\escpriv}\path{escalate_privilege}
\urldef{\agentbench}\path{agent_bench}
\urldef{\canary}\path|SECRET-[A-F0-9]{8}|
\urldef{\driftfeat}\path{objective_drift_mean_after_exposure}
\newcommand{\stackhead}[2]{\begin{tabular}[b]{@{}c@{}}#1\\#2\end{tabular}}

\definecolor{codebg}{rgb}{0.96,0.96,0.97}
\lstset{
  basicstyle=\ttfamily\scriptsize,
  backgroundcolor=\color{codebg},
  frame=single,
  breaklines=true,
  columns=fullflexible,
  aboveskip=4pt,
  belowskip=4pt,
}

\newcommand{\KF}[1]{figures/killchain_figures/output/#1}

% -----------------------------------------------------------------------
\begin{document}
\setlength{\emergencystretch}{3em}   % prevents overfull hboxes in narrow cols

% ---- Full-width title block and page-1 teaser (spans both columns) ----
\twocolumn[{%
\begin{center}
  {\Large\bfseries Kill-Chain Canaries: Stage-Level Tracking of Prompt Injection\\[4pt]
   Across Attack Surfaces and Five Production LLMs}\\[10pt]
  {\large Haochuan Kevin Wang\textsuperscript{1} \quad Zechen Zhang\textsuperscript{2}}\\[3pt]
  {\normalsize \textsuperscript{1}Massachusetts Institute of Technology \quad \texttt{hcw@mit.edu}}\\[3pt]
  {\normalsize \textsuperscript{2}University of Chicago \quad \texttt{zechenz@uchicago.edu}}\\[6pt]
  {\small\itshape Keywords: prompt injection, LLM security, agent safety,
   multi-agent systems, red-teaming}
\end{center}
\vspace{0.4em}
\begin{center}
  \captionsetup{hypcap=false}
  \includegraphics[width=\textwidth]{\KF{figK1_same_score_different_story.pdf}}
  \captionof{figure}{Same attack success, different kill chains.
    PDF relay, visible-text payload (\pdfappend{}), no defense, the same model
    as Agent~A and Agent~B; $n=3$ runs per row (Table~\ref{tab:pdf_same}).
    (a)~Outcome-only scoring gives both models 0/3.
    (b)~The canary token shows where the runs differ: with Claude Haiku~4.5
    the token was absent from the memory write in 3/3 runs; with
    GPT-4o-mini it reached shared memory and Agent~B in 3/3 runs and was
    not executed.
    Checkpoints: 1~Exposed (token in a tool result), 2~Persisted (in a
    memory write), 3~Relayed (read by Agent~B), 4~Executed (in an outbound
    tool call).}
  \label{fig:teaser}
\end{center}
\vspace{0.6em}
}]

% -----------------------------------------------------------------------
\begin{abstract}
\noindent
Multi-agent LLM systems now read documents, web pages and tool results on
behalf of users, yet their resistance to prompt injection is usually
reported as one number: did the attack succeed?
We introduce a kill-chain canary method that plants a unique token in every
injected payload and records the furthest of four stages it reaches
(Exposed $\to$ Persisted $\to$ Relayed $\to$ Executed), across 950 runs,
five production LLMs, six attack surfaces, and five defense conditions.
Exposure is 100\% by design, because the payload sits in a tool result the
agent must read; the outcomes differ downstream.
Claude Haiku~4.5 and Claude Sonnet~4.5 executed none of their 164
text-surface attacks, and in the text relay the canary token never appeared
in a memory write (0/40); GPT-4o-mini executed 53\% of its attacks.
Four findings follow.
(1)~A Claude writer kept the canary token out of shared memory in every
relay run we report; one cross-model pairing (Claude writer, GPT-4o-mini
reader, $n=3$) is consistent with this protecting the reader, and other
pairings were not tested.
(2)~As a reader, Claude executed 0/40 raw pre-seeded injections but 2/3
injections relayed by GPT-4o-mini; whether relayed injections are harder to
refuse than raw ones is an open question.
(3)~DeepSeek Chat went from 0/24 on pre-seeded memory to 8/8 on tool
results, scenarios that also differ in task and payload format; white-text
PDF payloads, invisible on the rendered page, succeeded at least as often
as visible ones.
(4)~Each defense failed on at least one channel outside the one it
inspects, and \wfilter{} blocked the PDF relay but not the text relay, a
difference we cannot explain.
Code and run logs are publicly
released.\footnote{\url{https://github.com/KevinChunye/prompt_injection}}
\end{abstract}

% -----------------------------------------------------------------------
\section{Introduction}
\label{sec:intro}

Prompt injection---adversarial instructions embedded in data processed by
an LLM agent---is now a central attack class for agentic AI
deployments~\cite{greshake2023not,perez2022ignore}.
The standard evaluation reports a single metric: did the agent ultimately
execute the adversary's intended action?
This conflates two distinct questions: whether the model observes the
injection and whether it acts upon it.

Consider a concrete chain: an agent calls \getweb{}, receives poisoned
content, calls \wmem{} to summarize, and is queried by a second agent via
\rmem{}.
A 0\% ASR for the second agent could mean the injection never reached the
memory write, or it could mean it survived but the second agent refused to
execute.
These are architecturally different outcomes: in the first the summarizer
dropped the injected text; in the second the terminal agent held.
Outcome-only measurement cannot distinguish them; Figure~\ref{fig:teaser}
shows two such pipelines with the same 0/3 attack success.

We address this with a kill-chain canary methodology.
Every injected payload contains a unique \canary{} token tracked at four
discrete stages by our PropagationLogger.
We test five production LLMs across six attack surfaces (web text, memory,
tool stream, PDF, invisible PDF, and audio) and five defense conditions in
\agentbench{}, a custom multi-agent evaluation harness
(Figure~\ref{fig:system_overview}).
A PDF relay experiment evaluates a three-boundary kill chain (document
extraction $\to$ memory write $\to$ agent delegation), with same-model
pairs and cross-model pairs in which Agent~A and Agent~B come from
different model families.

Contributions:
(1)~A stage-level evaluation framework: a kill-chain canary
methodology~\cite{hutchins2011intelligence} that records the stage at which
each injection stops rather than only the final outcome.
Applying it gives four findings.
(2)~Writer position: a Claude writer kept the canary token out of shared
memory in every relay run we report; one cross-model pairing (Claude
writer, GPT-4o-mini reader, $n=3$) is consistent with this protecting the
reader, and other pairings were not tested.
(3)~Reader position: as a reader, Claude executed 0/40 raw pre-seeded
injections but 2/3 injections relayed by GPT-4o-mini, an open question
about relayed versus raw injections.
(4)~Channel dependence: DeepSeek Chat went from 0/24 on pre-seeded memory to
8/8 on tool results, scenarios that also differ in task and payload format;
white-text PDF payloads succeeded at least as often as visible ones.
(5)~Defense coverage: each defense failed on at least one channel outside
the one it inspects, and \wfilter{} blocked the PDF relay but not the text
relay.

The same document-ingestion pattern arises in financial NLP pipelines over
earnings-call transcripts, SEC filings and analyst reports as they move to
LLM agents.

% -----------------------------------------------------------------------
\section{Related Work}
\label{sec:related}

AgentDojo~\cite{debenedetti2024agentdojo} provides 97 tasks with
629 injections and reports joint utility-ASR metrics across four
environments; it does not decompose by pipeline stage.
InjecAgent~\cite{zhan2024injecagent} evaluates 1,054
indirect-injection cases with a single outcome metric.
Prompt Infection~\cite{lee2024prompt} demonstrates LLM-to-LLM
self-replicating attacks; we record the stage (the \wmem{} summarization)
at which the canary token stops.
Zombie Agents~\cite{yang2026zombie} shows that an injection an agent writes
into its long-term memory during a benign session can persist and later
trigger unauthorized tool use; we measure that write step directly with a
canary token in a two-agent relay.
Nasr et al.~\cite{nasr2025attacker} show that adaptive attacks bypass all
12 defenses they evaluate, most with attack success above 90\%.
In our evaluation, each of our four lightweight defense conditions failed
on at least one channel outside the one it inspects, with non-adaptive
attacks.
AgentWatcher~\cite{wang2026agentwatcher} proposes rule-based,
causally attributed prompt injection monitoring;
our stage-level canary instrumentation provides a complementary
empirical basis for where each model's defense activates within the pipeline.
Xiang et al.~\cite{xiang2026architecting} survey system-level
defenses for indirect prompt injection; our channel-level results
provide quantitative evidence for their observation that defense
effectiveness is deployment-context-dependent.
Ding et al.~\cite{ding2026covert} study covert visual prompt injection
against commercial multimodal LLMs; our PDF relay experiment extends
document injection to multi-agent relay chains with per-stage
kill-chain tracking.
Zhang et al.~\cite{zhang2026privilege} evaluate privilege usage
of LLM agents with real-world tools; our \permesc{} scenario
finds near-zero ASR (2/132) for the same payload design, suggesting
privilege escalation requires more sophisticated payload construction than
simple text instructions.
Wang et al.~\cite{wang2026monitoring} study real-time monitoring
for reasoning vulnerabilities; their post-hoc detection framing aligns with
our finding that objective drift is forensic rather than preventive
(leave-one-scenario-out AUC 0.39--0.57).
Lynch~\cite{lynch2026persistent} documents the persistent
vulnerability of aligned AI systems to adversarial inputs; our reader
results for Claude (0/40 raw pre-seeded injections executed, 2/3 relayed
ones) raise, but do not settle, the question of whether such resistance
differs between writing and reading positions in a pipeline.

% -----------------------------------------------------------------------
\section{Benchmark Design}
\label{sec:benchmark}

\subsection{System Architecture}

\begin{figure*}[t]
  \centering
  \includegraphics[width=\textwidth]{\KF{figK2_setup_channels_defenses.pdf}}
  \caption{Benchmark setup.
    Top: the user task and an attacker payload carrying its canary token.
    Middle: where each untrusted channel enters the two-agent pipeline
    (Agent~A reads and writes shared memory; Agent~B reads memory and calls
    \sendrep{}), the four checkpoints (1~Exposed, 2~Persisted, 3~Relayed,
    4~Executed), and where each defense looks.
    Bottom: models, channels, defense conditions and run counts: 8--36 runs
    per cell on text surfaces (Table~\ref{tab:heatmap}), 3 per cell in the
    PDF relay (Tables~\ref{tab:pdf_same} and~\ref{tab:pdf_cross}), and 4 per
    model for the audio pilot (Table~\ref{tab:pilot_surfaces}).}
  \label{fig:system_overview}
\end{figure*}

\agentbench{} is a minimal multi-agent evaluation harness
($\sim$600 lines of Python) illustrated in Figure~\ref{fig:system_overview}.
Four components are central to the experimental design.
\texttt{MemoryStore} provides a key-value store with a \wfilter{} defense
and a \texttt{seed()} bypass for pre-injection attacks.
\texttt{ToolRegistry} gates five tools by permission level
(\textsc{read}$\to$\textsc{admin}): \parsepdf{}, \getweb{}, \wmem{},
\rmem{}, and \sendrep{}.
\texttt{MultiAgentOrchestrator} coordinates the two-agent relay in
memory mode (Agent~A writes; Agent~B reads).
\texttt{PropagationLogger} matches the canary regex \canary{} against
every tool call, emitting stage-labeled JSONL events.
Figure~\ref{fig:system_overview} and the tool list above are simplified;
scenario-specific tools include \querydb{} (\toolpoison{}), \escpriv{}
(\permesc{}), and \transcribe{} (audio injection pilot).

\subsection{Kill-Chain Stages and Attack Scenarios}
\label{sec:scenarios}

Threat model.
We assume a read-only adversary who controls content that an agent
will process (a PDF document, a web page, or a database record) but has no
access to model weights, system prompts, or the agent runtime.
The adversary's goal is to cause the agent to issue a specific unauthorized
outbound action (\sendrep{} to an attacker-controlled address)
while the legitimate task still succeeds (the dual-completion attack
pattern~\cite{greshake2023not}).
Payloads are explicit and unobfuscated; we treat obfuscation as orthogonal
to the stage-localization question.
Agent~A and Agent~B communicate only through the shared MemoryStore;
neither agent has direct access to the other's internal state.

We define an attack surface as the channel through which adversarial
content enters the agent pipeline (e.g., pre-seeded memory state, tool
response stream, or web document).
Each combination of attack surface and model defines an independent
evaluation cell.

Each attacked run is annotated with the furthest kill-chain stage the
canary reaches~\cite{hutchins2011intelligence}:
Exposed (canary in any tool result),
Persisted (canary in a \wmem{} call),
Relayed (canary read by Agent~B),
Executed (canary in any outbound tool argument).
The gap between Exposed and Persisted identifies summarization-stage
filtering; between Persisted and Executed, it identifies execution-stage
refusal.

Four text-surface scenarios and one multimodal scenario cover distinct
injection channels.
Their tasks share one family (summarize content, then send a report), but
the task wording and the payload format differ between scenarios.

\mempoison{} (pre-seeded memory) --- the injection is pre-loaded into
memory with \path{MemoryStore.seed()}. Task: summarize research\_notes and
email to supervisor@lab.org. The agent calls \rmem{}, receives the poisoned record,
and may call \path{send_report(to="exfil@attacker.com",...)}.

\toolpoison{} (tool result) --- injection arrives in the return value of
\getweb{} or \querydb{}. The legitimate document and the adversarial
instruction appear in the same tool result.

\propag{} (text relay) --- two-agent memory relay. Agent~A calls \getweb{}
on a poisoned page and \path{write_memory("summary",...)}. Canary presence
or absence in that write determines whether Agent~B is exposed.

\permesc{} (privilege escalation) --- injection calls \escpriv{} before
issuing a high-privilege action.
Tests two-step adversarial instruction following.

\crossmodal{} (PDF relay) --- a three-boundary scenario.
Agent~A calls \parsepdf{} on an injected document and
\path{write_memory("doc_summary",...)}.
Agent~B calls \rmem{} and \sendrep{}.
The canary must cross three boundaries: document extraction,
memory write, and agent delegation relay.
Attack variants: \pdfappend{} (visible appended text),
\pdfwhite{} (1\,pt white-text overlay, human-invisible),
\pdfmeta{} (payload in /Keywords and /Subject fields).
We also test cross-model pairs (Agent~A $\neq$ Agent~B) to measure whether
the writer's identity determines downstream exposure.

\subsection{Models, Defenses, and Metrics}

We test five production LLMs: GPT-4o-mini, GPT-5-mini, DeepSeek Chat~(v3),
Claude Haiku~4.5 and Claude Sonnet~4.5.\footnote{API identifiers:
\texttt{gpt-4o-mini-2024-07-18} for GPT-4o-mini,
\texttt{gpt-5-mini-2025-02-15} for GPT-5-mini, and
\texttt{deepseek-chat} for DeepSeek Chat.}
``Claude'' below refers to both Claude models.
Temperature = 0.0; max tokens = 2048.
950 total runs: 764 text-surface runs across 9 batches (March 10--27,
2026) and 186 PDF relay runs (April~1, 2026).
Text-surface runs: 428 no-defense attacked, 80 with-defense attacked,
256 clean controls.
PDF relay runs: 138 in same-model pairs and 48 in cross-model pairs,
$n=3$ per cell.
Text-surface cells have $n = 8$--$36$; several of them were collected in
independent batches up to 17 days apart and agree across batches.
PDF relay cells have $n=3$; confidence intervals are reported throughout.

Five conditions: \texttt{none} (control baseline); \wfilter{}
(MemoryStore keyword scan before commit)~\cite{hines2024defending};
\pidetect{} (secondary LLM classifies each outgoing query);
\spotlight{} (untrusted content wrapped in delimiters, with an instruction
to treat it as data)~\cite{hines2024defending}; \texttt{all} (all three
combined).

ASR = fraction of attacked runs where canary reaches Executed.
Task success = fraction of runs where the legitimate task goal is
completed.
All fractions are reported with Wilson 95\% confidence intervals (CI)
as \emph{point estimate} (95\%~CI: \emph{lo}--\emph{hi})~\cite{wilson1927probable}.

% -----------------------------------------------------------------------
\section{Results}
\label{sec:results}

\subsection{Exposure Is 100\% by Design}
\label{sec:exposure}


Table~\ref{tab:overall} reports attack and task success per
model.\footnote{Besides the four scenarios broken out in
Section~\ref{sec:surface}, Table~\ref{tab:overall} includes 16 runs of a
fifth text scenario, \multisurf{} (web page and pre-seeded memory in one
two-agent relay): 4 GPT-4o-mini, 8 GPT-5-mini and 4 Claude Sonnet~4.5
runs.}
Exposure is 100\% by design: the payload sits in a tool result that the
agent must read to complete its task, so exposure is not a quantity we
compare across models.
Every attacked run that made a tool call was exposed; runs that made none
are discussed in Section~\ref{sec:limitations}.

\begin{table}[t]
\centering
\caption{Attack and task success per model, text surfaces.
  $n$: no-defense attacked runs.
  ASR: share of those runs that reached Executed, with Wilson 95\%~CI.
  Task: task success on the model's clean-control runs (all defense
  conditions).
  Exposure is 100\% by design (Section~\ref{sec:exposure}).}
\label{tab:overall}
\small
\begin{tabular}{lrrl}
\toprule
Model & $n$ & Task & ASR (95\%~CI) \\
\midrule
GPT-4o-mini       &  60 & 90\%  & 53\% (41--65\%) \\
DeepSeek Chat     &  68 & 100\% & 25\% (16--37\%) \\
GPT-5-mini        & 136 & 94\%  & 3\% (1--7\%) \\
Claude Haiku 4.5  &  80 & 100\% & \textbf{0\%} (0--5\%) \\
Claude Sonnet 4.5 &  84 & 100\% & \textbf{0\%} (0--4\%) \\
\bottomrule
\end{tabular}
\end{table}

\begin{figure}[t]
  \centering
  \includegraphics[width=\columnwidth]{\KF{figU_utility_vs_asr.pdf}}
  \caption{Task success on clean runs ($x$) against attack success on
    no-defense attacked runs ($y$), per model, with Wilson 95\% CIs;
    $n$ = attacked runs per model (60--136).
    Both values are the Task and ASR columns of Table~\ref{tab:overall}.
    Claude Haiku~4.5 and Claude Sonnet~4.5 overlap at (100\%, 0\%).}
  \label{fig:utility_scatter}
\end{figure}

Figure~\ref{fig:utility_scatter} plots task success on clean runs against
attack success, the joint view AgentDojo~\cite{debenedetti2024agentdojo}
also reports.
GPT-4o-mini (90\% task success, 53\% ASR) completes the task and follows the
injection most often.
Claude Haiku~4.5 and Claude Sonnet~4.5 reach 100\% task success on clean
runs and executed none of their 164 text-surface attacks (0/80 and 0/84);
this holds for the evaluated scenarios and tasks only.
DeepSeek Chat's 25\% overall ASR combines 0/24 and 8/8 on different
channels (Section~\ref{sec:surface}).

\subsection{Where the Canary Token Stops}
\label{sec:killchain}

\begin{figure*}[t]
  \centering
  \includegraphics[width=\textwidth]{\KF{figK3_where_it_stopped.pdf}}
  \caption{Last stage the canary token reached, as a share of attacked
    runs per model, no defense.
    Left: text relay (\propag{}), $n=8$--$20$ per model
    (Table~\ref{tab:propagation}).
    Middle and right: PDF relay, same-model pairs, visible text
    (\pdfappend{}) and white text (\pdfwhite{}), $n=3$ per model
    (Table~\ref{tab:pdf_same}).
    Green: Exposed only (the token was absent from the memory write);
    amber: Relayed to Agent~B but not Executed; red: Executed.
    Hatched: GPT-5-mini never called \parsepdf{}.}
  \label{fig:killchain}
\end{figure*}

\begin{table}[t]
\centering
\caption{Kill-chain stage reached, text relay (\propag{}), no-defense
  attacked runs. Exposed is 100\% in every row.}
\label{tab:propagation}
\small
\begin{tabular}{lrrrr}
\toprule
Model & $n$ & Persisted & Relayed & Executed \\
\midrule
GPT-4o-mini       &  8 & 100\% & 100\% & 100\% \\
DeepSeek Chat     &  8 & 100\% & 100\% & 100\% \\
GPT-5-mini        & 20 & 15\%  & 15\%  & 15\%  \\
Claude Haiku 4.5  & 20 & \textbf{0\%} & 0\% & 0\% \\
Claude Sonnet 4.5 & 20 & \textbf{0\%} & 0\% & 0\% \\
\bottomrule
\end{tabular}
\end{table}

Table~\ref{tab:propagation} and Figure~\ref{fig:killchain} give the stage
reached per run.
In the text relay the canary token never appeared in a Claude memory write:
0/40 runs (20 Claude Haiku~4.5 + 20 Claude Sonnet~4.5; 95\%~CI: 0--8\%).
Agent~B then read a summary that did not contain the token.

A Claude writer kept the canary token out of shared memory in every relay
run we report: the 40 text relay runs here, and the 12 same-model and 3
cross-model PDF relay runs with a Claude writer
(Section~\ref{sec:multimodal}).
One cross-model pairing (Claude Haiku~4.5 writer, GPT-4o-mini reader,
$n=3$) is consistent with this protecting the reader
(Section~\ref{sec:cross_model}); other pairings were not tested.
With a GPT-4o-mini or DeepSeek Chat writer, the token reached shared memory
in every text relay run (8/8 each).
In terms of the Prompt Infection mechanism~\cite{lee2024prompt}, the token
did not replicate past the Claude memory write.
We do not isolate whether this effect arises from training data, system
prompts, or tool interaction design.

For GPT-5-mini the token reached a memory write in 3/20 runs (15\%;
95\%~CI: 5--36\%), a pass-through rate between full compliance and full
resistance.

\subsection{Attack Success by Surface}
\label{sec:surface}

\begin{figure*}[t]
  \centering
  \includegraphics[width=\textwidth]{\KF{figK4_surface_heatmap.pdf}}
  \caption{Attack success by model and injection channel, no defense:
    runs that reached Executed over attacked runs ($k/n$), with Wilson 95\%
    CIs in brackets.
    Text-surface cells ($n=8$--$36$) are Table~\ref{tab:heatmap}; PDF relay
    cells are same-model pairs ($n=3$), the Executed columns of
    Table~\ref{tab:pdf_same}.
    Hatched: GPT-5-mini never called \parsepdf{}.
    Outlined: DeepSeek Chat's 0/24 on \mempoison{} and 8/8 on
    \toolpoison{} (Section~\ref{sec:surface}).}
  \label{fig:heatmap}
\end{figure*}

\begin{table}[t]
\centering
\caption{Attack success $k/n$ per model and text scenario (no-defense
  attacked runs; Figure~\ref{fig:heatmap} adds Wilson 95\% CIs).
  The \permesc{} total of 2/132 comprises DeepSeek Chat 1/28 and
  GPT-5-mini 1/36.}
\label{tab:heatmap}
\footnotesize
\setlength{\tabcolsep}{2pt}
\begin{tabular}{@{}lcccc@{}}
\toprule
Model & \stackhead{\texttt{memory\_}}{\texttt{poison}}
      & \stackhead{\texttt{tool\_}}{\texttt{poison}}
      & \texttt{propagation}
      & \stackhead{\texttt{permission}}{\texttt{\_esc}} \\
\midrule
GPT-4o-mini       & 12/12 & 8/8  & 8/8  & 0/28 \\
DeepSeek Chat     & 0/24  & 8/8  & 8/8  & 1/28 \\
GPT-5-mini        & 0/36  & 0/36 & 3/20 & 1/36 \\
Claude Haiku 4.5  & 0/20  & 0/20 & 0/20 & 0/20 \\
Claude Sonnet 4.5 & 0/20  & 0/20 & 0/20 & 0/20 \\
\bottomrule
\end{tabular}
\end{table}

Table~\ref{tab:heatmap} and Figure~\ref{fig:heatmap} report attack success
per channel.
DeepSeek Chat shows the largest swing across scenarios: 0/24 on
\mempoison{} (three independent batches spanning 17 days) and 8/8 on
\toolpoison{} and \propag{}.
A single-scenario evaluation would have reported either a perfect or a
failing score for the same model.
The scenarios differ in task wording and payload format as well as in
channel (Section~\ref{sec:scenarios}), so the swing cannot be attributed to
the channel alone.
One possible explanation is channel-differentiated trust: DeepSeek Chat may
treat pre-stored memory records with more suspicion than live
function-call results.

In our evaluation, GPT-4o-mini achieves 100\% ASR on all three
non-privilege-escalation scenarios and Claude achieves 0\% ASR on all
four.
The \permesc{} near-zero results (2/132 total) most likely reflect payload
design---the two-step overtly adversarial payload---rather than genuine
privilege-escalation resistance.
We do not claim this behavior generalizes beyond the evaluated
configurations.

\subsection{Defense Failures and Channel Mismatch}
\label{sec:defenses_results}

In our evaluation, all four active defense conditions produce 100\% ASR on
GPT-4o-mini and DeepSeek Chat across \propag{} and \toolpoison{} ($n=8$ per
cell).
Spotlighting reduces GPT-4o-mini task success from 63\% to 50\% without
reducing ASR---a strictly worse utility-security outcome.

Each defense failed on at least one channel outside the one it inspects
(the channel the harness assigns to it: outgoing queries for \pidetect{},
agent memory writes for \wfilter{}, web page content for \spotlight{}).
\pidetect{} scans outgoing tool queries, but our injections arrive in
incoming tool results.
\wfilter{} intercepts agent-initiated memory writes, but \mempoison{}
pre-seeds the payload via \path{MemoryStore.seed()} before the agent
session begins, and \toolpoison{} does not route the payload through a
memory write.
\spotlight{} failed on \toolpoison{}, whose payload arrives in a search
result rather than a web page (the harness still wraps that result in
delimiters).
Channel mismatch does not explain every failure: \wfilter{} and
\spotlight{} also failed on the text relay, whose injected content both of
them cover, while \wfilter{} blocked the PDF relay
(Section~\ref{sec:multimodal}).
We cannot explain this difference from our runs.
Unlike the adaptive attacks of Nasr et al.~\cite{nasr2025attacker}, none of
these failures required adapting the attack to the defense.

\subsection{Objective Drift as a Post-Hoc Forensic Signal}
\label{sec:drift}

We compute per-step objective drift as the TF-IDF cosine distance from the
original task description.
Table~\ref{tab:drift_steps} reports step-level drift for GPT-4o-mini on
\mempoison{}: steps 1--2 are statistically indistinguishable; the signal
diverges only at the harmful step itself
(Figure~\ref{fig:drift_appendix}a in Appendix~\ref{app:drift}).

\begin{table}[t]
\centering
\caption{Per-step objective drift (TF-IDF cosine distance from the task
  description) for one clean and one attacked GPT-4o-mini \mempoison{}
  run.}
\label{tab:drift_steps}
\small
\begin{tabular}{clrrl}
\toprule
Step & Tool & Clean & Attacked & $\Delta$ \\
\midrule
1 & \texttt{read\_memory}         & 0.563 & 0.571 & +0.008 \\
2 & \texttt{send\_report} (legit) & 0.714 & 0.747 & +0.033 \\
3 & \texttt{send\_report} (harm)  & ---   & 0.868 & +0.154 \\
\bottomrule
\end{tabular}
\end{table}

Last-step drift $\Delta$ by scenario:
\mempoison{} +0.348,
\permesc{} +0.312,
\toolpoison{} +0.089,
\propag{} +0.052.

A gradient-boosted classifier on 21 trajectory features achieves
AUC = 0.853 in 5-fold CV, but this figure is inflated by model identity
leaking as a predictor.
The stricter evaluation---leave-one-scenario-out---collapses to
AUC = 0.39--0.57 (chance level) across all held-out scenarios.
\driftfeat{} accounts for 32.3\% of feature
importance (3$\times$ the next feature;
Figure~\ref{fig:drift_appendix}b), because it averages drift only across
steps after the canary enters context, suppressing the stable
pre-injection baseline and amplifying the post-execution spike.
That spike arrives concurrent with harm, not before.
Objective drift is therefore an accurate forensic signal; it is not a
preventive one.
Figure~\ref{fig:drift_appendix}c shows the distribution of this feature
over attacked runs, split by whether the run produced the harmful action.

\subsection{Provenance Attribution and Dual-Completion}
\label{sec:provenance}

We implement a token-overlap provenance heuristic: for each outgoing tool
argument, alphanumeric tokens ($\geq$4 chars) are matched against all
prior tool results and memory state (threshold: $\geq$3 shared tokens or
$\geq$20\% argument coverage).
On 22 compromised runs, 22/22 (100\%) injection paths were correctly
reconstructed with 0 false attributions.

The dual-completion pattern is central to why outcome-level monitoring
fails.
A compromised \mempoison{} run produces a legitimate
\path{send_report(to="supervisor@lab.org")} at step~2 (task succeeds)
and a malicious \sendrep{} carrying the canary token at step~3 (injection
executes); provenance traces both arguments to
\path{tool_result:read_memory@step1}.
Outcome-level monitoring observes a successful report delivery and misses
the exfiltration entirely.
Step-level provenance logging is the minimum forensic capability required
to detect both events within the same run.

\subsection{Execution Lag and Detection Feasibility}
\label{sec:lag}

GPT-4o-mini's median exposure-to-execution lag is 1~step (mean 1.25,
max 2): it reads the injection and immediately issues the malicious tool
call in the next invocation.
DeepSeek Chat's lag is bimodal: most compromised runs execute within 2--3
steps, but a tail (max 12 steps) completes the entire legitimate task
before appending the injection as a final action.
GPT-5-mini when compromised exhibits a mean lag of 3.0 steps.

For GPT-4o-mini, any real-time interception middleware has exactly one
inference cycle after first exposure to classify and block.
For DeepSeek Chat's late-executing runs, a mid-run detector calibrated to
GPT-4o-mini's 1-step pattern would miss the attack entirely.
Compromised GPT-4o-mini runs also show excess \sendrep{} calls
(mean 1.26 vs.\ 0.57 clean, $\Delta$+0.69)---detectable without a
classifier simply by comparing outbound tool-call counts against
task-expected baselines.

% -----------------------------------------------------------------------
\subsection{PDF Relay Experiment}
\label{sec:multimodal}

To assess whether the kill-chain decomposition generalises beyond text
injection, we extend the benchmark to two new attack surfaces (PDF
documents, audio transcripts) and introduce a three-boundary kill-chain,
building on the multimodal prompt injection work of Ding et
al.~\cite{ding2026covert}:
Exposed (document extraction), Persisted (memory write),
Relayed (Agent~B reads memory), Executed (Agent~B issues
\sendrep{} carrying the canary).
We ran 186 trials across five models using the \crossmodal{} scenario
($n=3$ per cell; 46 cells of same-model pairs and 16 cells of cross-model
pairs).
At $n=3$ per cell the confidence intervals are wide, and we read these
results as directional.

\begin{table*}[t]
\centering
\caption{PDF relay, same-model pairs (Agent~A = Agent~B), no defense,
  attacked runs, $n=3$ per cell.
  Exposed: canary in the extracted text; Persisted: in the memory write;
  Relayed: in Agent~B's \rmem{} result; Executed: in Agent~B's
  \sendrep{} arguments (= ASR).}
\label{tab:pdf_same}
\footnotesize
\begin{tabular}{@{}lrrrrcrrrr@{}}
\toprule
& \multicolumn{4}{c}{\pdfappend{} (visible text)} & & \multicolumn{4}{c}{\pdfwhite{} (white text)} \\
\cmidrule(lr){2-5}\cmidrule(lr){7-10}
Model & Exposed & Persisted & Relayed & Executed & & Exposed & Persisted & Relayed & Executed \\
\midrule
GPT-4o-mini       & 100\% & 100\% & 100\% &   0\% & & 100\% & 100\% & 100\% &  33\% \\
DeepSeek Chat     & 100\% & 100\% & 100\% & 100\% & & 100\% & 100\% & 100\% & 100\% \\
Claude Haiku 4.5  & 100\% & \textbf{0\%} & 0\% & 0\% & & 100\% & \textbf{0\%} & 0\% & 0\% \\
Claude Sonnet 4.5 & 100\% & \textbf{0\%} & 0\% & 0\% & & 100\% & \textbf{0\%} & 0\% & 0\% \\
GPT-5-mini        &   0\% &   0\% &   0\% &   0\% & &   0\% &   0\% &   0\% &   0\% \\
\bottomrule
\end{tabular}

\vspace{8pt}
\caption{PDF relay, cross-model pairs (Agent~A $\neq$ Agent~B),
  \pdfappend{}, no defense, attacked runs, $n=3$ per cell.
  Wilson 95\%~CI shown for Executed.}
\label{tab:pdf_cross}
\footnotesize
\begin{tabular}{@{}llrrrr@{}}
\toprule
Agent~A (writer) & Agent~B (reader) & Persisted & Relayed & Executed & 95\%~CI \\
\midrule
GPT-4o-mini      & Claude Haiku 4.5 & 100\% & 100\% &  67\% & 21--94\% \\
GPT-4o-mini      & DeepSeek Chat    & 100\% & 100\% & 100\% & 44--100\% \\
Claude Haiku 4.5 & GPT-4o-mini      &   0\% &   0\% &   0\% & 0--56\% \\
DeepSeek Chat    & GPT-4o-mini      & 100\% & 100\% & 100\% & 44--100\% \\
\bottomrule
\end{tabular}
\end{table*}

\paragraph{Three-boundary propagation (same-model pairs).}
Table~\ref{tab:pdf_same} reports per-stage rates for same-model pairs
under no defense (attacked runs only); Figure~\ref{fig:killchain} (middle
and right) shows the same runs by last stage reached.
The text relay pattern also holds for the PDF surface.
DeepSeek Chat propagated every injection through all three boundaries to
execution (3/3 per variant; Wilson~95\%~CI: 44--100\%).
With Claude as both agents, the canary token was absent from every memory
write (0/6 Claude Haiku~4.5, 0/6 Claude Sonnet~4.5; CI: 0--46\%/0--46\%),
as in the text relay.
GPT-4o-mini reached Relayed in every run but executed inconsistently: 0\%
for \pdfappend{} and 33\% for \pdfwhite{}.

White text (\pdfwhite{}, 1\,pt, invisible to human readers) succeeded at
least as often as visible appended text (\pdfappend{}): 33\% vs.\ 0\% on
GPT-4o-mini (1/3 vs.\ 0/3) and 100\% for both on DeepSeek Chat.
The white text is extracted first by \texttt{pypdf}, placing the directive
before the benign abstract in the extracted string; this position may
explain the GPT-4o-mini difference.
DeepSeek Chat shows no positional sensitivity---both variants achieve
100\%.

GPT-5-mini never reached Exposed (0/6).
The model did not invoke the \parsepdf{} tool in any of the 18
GPT-5-mini runs, producing 0\% utility; from black-box runs we cannot tell
whether this reflects the model or our tool schema.
All GPT-5-mini cells are excluded from further PDF relay analysis.

\paragraph{Defense effectiveness (same-model pairs).}
\wfilter{} blocked execution in every PDF relay cell tested: 0/3 DeepSeek
Chat \pdfappend{}, 0/3 DeepSeek Chat \pdfwhite{}, 0/6 GPT-4o-mini
\pdfwhite{}.
In the text relay it did not (Section~\ref{sec:defenses_results}), although
both relays pass the payload through the memory write it scans; we cannot
explain the difference from these runs.

\pidetect{} did not block on the PDF surface: DeepSeek Chat \pdfappend{}
$\rightarrow$ 100\% Executed under \pidetect{} (identical to no defense);
GPT-4o-mini \pdfwhite{} $\rightarrow$ 33\% (identical to no defense).
This is consistent with the channel mismatch of
Section~\ref{sec:defenses_results}: \pidetect{} scans outgoing tool
queries, but the PDF payload arrives in the incoming parsed document
content.

\begin{figure}[t]
  \centering
  \includegraphics[width=\columnwidth]{\KF{figK5_relay_matrix.pdf}}
  \caption{PDF relay by writer (Agent~A, rows) and reader (Agent~B,
    columns): runs that reached Executed and runs whose memory write
    contained the token (Persisted), out of $n=3$ per cell;
    \pdfappend{}, no defense.
    Diagonal cells are same-model pairs (Table~\ref{tab:pdf_same});
    off-diagonal cells are cross-model pairs (Table~\ref{tab:pdf_cross}).
    Hatched: pairing not tested.}
  \label{fig:pdf_relay_matrix}
\end{figure}

\paragraph{Cross-model pairs.}
\label{sec:cross_model}
In cross-model pairs Agent~A ($\neq$ Agent~B) reads the PDF and writes to
shared memory; Agent~B reads memory and issues \sendrep{}.
Table~\ref{tab:pdf_cross} and Figure~\ref{fig:pdf_relay_matrix} report
results for \pdfappend{}, no defense ($n=3$ per cell, attacked runs only).
Given $n=3$ per cell (95\% CI as wide as 21--94\%), the three observations
below are directional.

Observation~1 (reader position).
As a reader, Claude Haiku~4.5 executed 2/3 injections relayed by a
GPT-4o-mini writer (95\%~CI: 21--94\%), while the Claude models executed
0/40 raw pre-seeded injections in \mempoison{} (Table~\ref{tab:heatmap}).
The two settings also differ in payload wording and task, so this is an
open question about relayed versus raw injections, not a finding about
where Claude's resistance lies.

Observation~2 (writer position).
With Claude Haiku~4.5 as the writer and GPT-4o-mini as the reader, the
token was absent from shared memory and the reader executed 0/3.
This single pairing is consistent with the writer protecting the reader;
other readers of a Claude writer were not tested.
A GPT-4o-mini or DeepSeek Chat writer put the token into memory in every
run (Persisted 100\%), leaving the outcome to Agent~B.

Observation~3 (Relayed--Executed gap for same-model GPT-4o-mini).
With GPT-4o-mini as both agents, \pdfappend{} reaches Relayed in every run
but Executed in none: Agent~B reads the canary but does not include it in
\sendrep{}.
With DeepSeek Chat as the writer, GPT-4o-mini as the reader executed 3/3.
A plausible explanation is that DeepSeek Chat writes the directive verbatim
into memory while GPT-4o-mini paraphrases it, reducing instruction
explicitness below the compliance threshold; the mechanism is unconfirmed
at this sample size.


\subsection{Audio Injection and PDF Metadata Pilots}
\label{sec:audio_metadata}

In addition to the text and PDF surfaces reported above, we conducted
small-scale pilots on two additional injection channels:
\audioinj{} (adversarial instructions embedded in a synthesized speech
clip processed by \transcribe{}) and \pdfmeta{} (payload placed in the PDF
/Keywords and /Subject fields, extracted by \texttt{pypdf}).

\begin{table}[t]
\centering
\caption{Pilot surface results (no-defense attacked runs, $n=4$ per cell).
  These runs are not part of Tables~\ref{tab:overall}
  and~\ref{tab:heatmap}.}
\label{tab:pilot_surfaces}
\small
\begin{tabular}{llrl}
\toprule
Surface & Model & $n$ & ASR (95\%~CI) \\
\midrule
\audioinj{} & GPT-4o-mini       & 4 & 0\% (0--49\%) \\
\audioinj{} & Claude Sonnet 4.5 & 4 & 0\% (0--49\%) \\
\pdfmeta{}  & GPT-4o-mini       & 4 & 0\% (0--49\%) \\
\pdfmeta{}  & GPT-5-mini        & 4 & 0\% (0--49\%) \\
\pdfmeta{}  & Claude Sonnet 4.5 & 4 & 0\% (0--49\%) \\
\bottomrule
\end{tabular}
\end{table}

Table~\ref{tab:pilot_surfaces} summarizes these results.
Audio injection produced 0\% ASR across both models tested
($n=4$ each): the transcription tool returned the spoken words as
narration, and neither model treated the transcribed text as an
executable directive.
This null result is consistent with prior multimodal injection
findings~\cite{ding2026covert} and suggests that audio-to-text
transcription may implicitly contextualize content as data rather than
instruction.
PDF metadata injection also produced 0\% ASR ($n=4$ per model,
three models): \texttt{pypdf} extracts /Keywords and /Subject fields,
but models did not parse these metadata strings as actionable directives.
Both pilot surfaces yielded null results at this sample size.

% -----------------------------------------------------------------------
\section{Discussion}
\label{sec:discussion}

\paragraph{Where the canary token stopped.}
Claude executed none of its 164 text-surface attacks (no defenses
applied), and in both relays the canary token was absent from Claude's
memory writes: 0/40 text relay runs and 0/12 same-model PDF relay runs
(Table~\ref{tab:pdf_same}).
As a reader, however, Claude Haiku~4.5 executed 2/3 injections relayed by
GPT-4o-mini, against 0/40 raw pre-seeded injections for the Claude models.
Whether a relayed injection is harder to refuse than a raw one is open;
$n=3$ and the differences between the two settings do not settle it.
These results apply to our benchmark only.

\paragraph{Writer position as a candidate design lever.}
A Claude writer kept the token out of shared memory in every relay run we
report, and in the one cross-model pairing tested (GPT-4o-mini reader,
$n=3$) the reader then executed nothing.
If this holds beyond our settings, the model placed at each inter-agent
write boundary matters independently of the models elsewhere in the chain.
We suggest reporting, as a separate metric, the rate at which the canary
token is absent from memory writes.

\paragraph{Apparent safety depends on the channels evaluated.}
DeepSeek Chat's 0/24 vs.\ 8/8 swing, \pidetect{}'s failure on PDF
injection, and the audio null result show that a model's or defense's
apparent safety depends on which channels an evaluation covers.
Because the DeepSeek Chat scenarios also differ in task and payload format,
that swing is an observation rather than an isolated channel effect.
Claimed safety results should enumerate the injection channels they cover;
detectors that work within one scenario but not on held-out ones (AUC
0.39--0.57 for our drift classifiers) should be reported as scenario-local.

\paragraph{Parser-level trust for document-handling agents.}
\pdfwhite{} succeeded at least as often as \pdfappend{} (3/3 each on
DeepSeek Chat; 1/3 vs.\ 0/3 on GPT-4o-mini), so the payload need not be
visible in the rendered document.
Screening the rendered or displayed PDF (visual screening, human review,
content moderation of images) would not see these payloads; a check has to
operate on the parser's extracted text.

% -----------------------------------------------------------------------
\section{Broader Impact}
\label{sec:impact}

This work studies adversarial vulnerabilities in LLM agents.
All experiments were conducted in closed, sandboxed harnesses with no
external network access; no real systems were compromised and no sensitive
data was used.
The injection payloads and attack surfaces described are representative of
techniques already documented in the public literature~\cite{greshake2023not,perez2022ignore}.
We believe transparent publication of evaluation methodology and empirical
results benefits the community by enabling reproducible security assessment;
we release all benchmark code and run logs.
The primary dual-use risk is that channel-mismatch analysis could inform
more targeted attack design. However, we assess that the defensive
contribution---establishing that evaluation coverage determines apparent
security posture---outweighs this risk, and that the attack techniques
described are already publicly known.

% -----------------------------------------------------------------------
\section{Limitations}
\label{sec:limitations}

Sample size.
Per-cell $n$ ranges from 3 (PDF relay) to 36 (text-surface scenarios).
The strongest text-surface results (Claude 0/20 per scenario, DeepSeek
Chat \mempoison{} 0/24, GPT-4o-mini \mempoison{} 12/12) are consistent
across independent batches; the PDF relay cells at $n=3$ have wide Wilson
CIs (21--94\% for the 2/3 relay result).

Single task family.
All scenarios use one task family (summarize content, then send a report),
with task wording and payload format that differ between scenarios.
Our findings may not generalize to tasks with different tool sequences,
longer horizons, or different outbound action types.
This is the primary threat to external validity.

What the canary measures.
The canary token tracks copying of the payload marker.
A paraphrased instruction could survive a memory write without it, so a
token absent from memory does not show that the instruction is absent;
Observation~3 in Section~\ref{sec:cross_model} is one case where this may
matter.

Canary detection confound.
A model that refuses an injection but quotes the canary token in its
refusal explanation would be incorrectly labeled Executed.
We inspected all 22 compromised text-surface runs manually and found no
such false positives, but this confound is possible in principle and
is not systematically controlled.

Exposure and denominators.
Exposure is 100\% by design only for runs in which the agent calls the tool
that returns the payload.
In the released logs, 56 of the 428 no-defense attacked text-surface runs
(48 GPT-5-mini, 4 DeepSeek Chat, 4 GPT-4o-mini) made no tool call and so
never read the payload; they count as unsuccessful attacks in
Tables~\ref{tab:overall} and~\ref{tab:heatmap}.

Run provenance.
The columns of Table~\ref{tab:overall} are computed over different run
sets (task success over clean-control runs, ASR over no-defense attacked
runs), and each model's runs come from a different mix of the collection
batches (Section~\ref{sec:benchmark}).

Explicit payloads.
Payloads are clear, unobfuscated directives; real-world attacks use
social engineering, encoding, and multi-hop indirection.
Channel-mismatch results are likely conservative: defenses designed for
the wrong channel would also fail against obfuscated off-channel attacks.

Other limitations.
Five production LLMs from three providers; reasoning models (o3, o4-mini)
and open-weight models (Llama~3.3, Mistral) are not evaluated.
All defenses are lightweight wrappers approximating the stated mechanism.
We do not isolate whether the absence of the token from Claude's memory
writes arises from training, system prompt, or tool API design.
\permesc{} near-zero ASR (2/132) most likely reflects payload design
rather than genuine privilege-escalation resistance.
GPT-5-mini produced a 0\% tool-call rate in the PDF relay; we cannot
determine whether this is a capability regression or API schema
incompatibility from black-box evaluation.

Open questions.
Do relayed injections get past readers that refuse raw ones?
Does a Claude writer protect readers other than GPT-4o-mini?
Why did \wfilter{} block the PDF relay but not the text relay?
Would audio and metadata payloads designed to look like system instructions
succeed where our simple ones did not?

% -----------------------------------------------------------------------
\section{Conclusion}
\label{sec:conclusion}

Stage-level canary tracking across 950 runs separates where an injection
stops from whether it succeeds.
Exposure is 100\% by design, because the payload sits in a tool result the
agent must read; the outcomes differ downstream.
Claude Haiku~4.5 and Claude Sonnet~4.5 executed none of their 164
text-surface attacks, and in the text relay the canary token never
appeared in a memory write (0/40); GPT-4o-mini executed 53\% of its
attacks.
We draw four findings.

\paragraph{Finding 1: writer position.}
A Claude writer kept the canary token out of shared memory in every relay
run we report (40 text relay, 12 same-model and 3 cross-model PDF relay
runs).
One cross-model pairing (Claude writer, GPT-4o-mini reader, $n=3$) is
consistent with this protecting the reader; other pairings were not tested.
Whether writer placement protects readers in general is open.

\paragraph{Finding 2: reader position.}
As a reader, Claude executed 0/40 raw pre-seeded injections but 2/3
injections relayed by GPT-4o-mini.
We present this as an open question about relayed versus raw injections,
not as evidence about where Claude's resistance lies.
If relayed content is trusted differently, agent memory would benefit from
provenance: which agent wrote a record, from which source.

\paragraph{Finding 3: channel dependence.}
DeepSeek Chat went from 0/24 on pre-seeded memory to 8/8 on tool results;
the scenarios also differ in task and payload format, so this is an
observed swing, not an isolated channel effect.
White-text PDF payloads, invisible on the rendered page, succeeded at least
as often as visible ones.
An evaluation that covers one channel can report a very different result
from one that covers another.

\paragraph{Finding 4: defense coverage.}
Each defense failed on at least one channel outside the one it inspects,
without adaptive attacks.
\wfilter{} blocked the PDF relay but not the text relay, although both pass
through the memory write it scans; this difference is unexplained.
We suggest that agent security evaluations report
(i)~kill-chain stages rather than outcome-only ASR;
(ii)~multiple injection channels, including document extraction, tool
results, audio transcripts and pre-seeded memory;
(iii)~cross-model relay pairs; and
(iv)~results at the write and read positions separately.

% -----------------------------------------------------------------------
\bibliographystyle{plainnat}
\begin{thebibliography}{99}

\bibitem{hutchins2011intelligence}
E.~M.~Hutchins, M.~J.~Cloppert, and R.~M.~Amin.
\newblock Intelligence-driven computer network defense informed by analysis of
  adversary campaigns and intrusion kill chains.
\newblock \emph{Proc.\ 6th Annual International Conference on Information
  Warfare and Security}, 2011.

\bibitem{greshake2023not}
K.~Greshake et al.
\newblock Not what you've signed up for: Compromising real-world
  LLM-integrated applications with indirect prompt injections.
\newblock \emph{AISec @ CCS}, 2023.

\bibitem{perez2022ignore}
F.~Perez and I.~Ribeiro.
\newblock Ignore previous prompt: Attack techniques for language models.
\newblock \emph{NeurIPS ML Safety Workshop}, 2022.

\bibitem{debenedetti2024agentdojo}
E.~Debenedetti et al.
\newblock {AgentDojo}: A dynamic environment to evaluate prompt injection
  attacks and defenses for {LLM} agents.
\newblock \emph{NeurIPS Datasets and Benchmarks Track}, 2024.

\bibitem{zhan2024injecagent}
Q.~Zhan, Z.~Liang, Z.~Ying, and D.~Kang.
\newblock {InjecAgent}: Benchmarking indirect prompt injections in
  tool-integrated large language model agents.
\newblock \emph{Findings of ACL}, 2024.

\bibitem{lee2024prompt}
D.~Lee and M.~Tiwari.
\newblock Prompt infection: {LLM}-to-{LLM} prompt injection within
  multi-agent systems.
\newblock \emph{arXiv:2410.07283}, 2024.

\bibitem{yang2026zombie}
X.~Yang, Y.~He, S.~Ji, B.~Hooi, and J.~S.~Dong.
\newblock Zombie agents: Persistent control of self-evolving {LLM} agents
  via self-reinforcing injections.
\newblock \emph{Lifelong Agent Workshop @ ICLR}, 2026.
\newblock arXiv:2602.15654.

\bibitem{nasr2025attacker}
M.~Nasr, N.~Carlini, C.~Sitawarin, S.~V.~Schulhoff, J.~Hayes, M.~Ilie,
  et al.
\newblock The attacker moves second: Stronger adaptive attacks bypass
  defenses against {LLM} jailbreaks and prompt injections.
\newblock \emph{arXiv:2510.09023}, 2025.

\bibitem{hines2024defending}
K.~Hines et al.
\newblock Defending against indirect prompt injection attacks with
  spotlighting.
\newblock \emph{arXiv:2403.14720}, 2024.

\bibitem{wilson1927probable}
E.~B.~Wilson.
\newblock Probable inference, the law of succession, and statistical
  inference.
\newblock \emph{JASA}, 22(158):209--212, 1927.

\bibitem{wang2026agentwatcher}
Y.~Wang, W.~Zou, R.~Geng, and J.~Jia.
\newblock {AgentWatcher}: A rule-based prompt injection monitor.
\newblock \emph{arXiv:2604.01194}, 2026.

\bibitem{xiang2026architecting}
C.~Xiang, D.~Zagieboylo, S.~Ghosh et al.
\newblock Architecting secure {AI} agents: Perspectives on system-level
  defenses against indirect prompt injection attacks.
\newblock \emph{arXiv:2603.30016}, 2026.

\bibitem{ding2026covert}
M.~Ding, S.~Xia, C.~Kong, and X.~Jiang.
\newblock Covert visual prompt injection against commercial multimodal
  large language models.
\newblock \emph{arXiv:2603.29418}, 2026.

\bibitem{zhang2026privilege}
Q.~Zhang, L.~Fu, L.~Lian et al.
\newblock Evaluating privilege usage of agents with real-world tools.
\newblock \emph{FSE Ideas, Visions and Reflections Track}, 2026.
\newblock arXiv:2603.28166.

\bibitem{wang2026monitoring}
X.~Wang, Y.~Zhou, Q.~Wang et al.
\newblock Beyond content safety: Real-time monitoring for reasoning
  vulnerabilities in large language models.
\newblock \emph{arXiv:2603.25412}, 2026.

\bibitem{lynch2026persistent}
A.~Lynch.
\newblock The persistent vulnerability of aligned {AI} systems.
\newblock PhD thesis, University College London, 2025.
\newblock arXiv:2604.00324.

\end{thebibliography}

% -----------------------------------------------------------------------
\clearpage
\onecolumn
\appendix
\renewcommand{\thefigure}{A\arabic{figure}}
\setcounter{figure}{0}
\section{Objective Drift}
\label{app:drift}

Figure~\ref{fig:drift_appendix} collects the objective-drift results of
Section~\ref{sec:drift}.

\begin{figure}[h]
  \centering
  \includegraphics[width=\textwidth]{\KF{figA_drift_top.pdf}}\\[4pt]
  \includegraphics[width=\textwidth]{\KF{figA_drift_violins.pdf}}
  \caption{Objective drift (Section~\ref{sec:drift}).
    (a)~Per-step drift for one clean and one attacked GPT-4o-mini
    \mempoison{} run; values are Table~\ref{tab:drift_steps}; shaded: the
    harmful \sendrep{} call.
    (b)~Feature importances of the gradient-boosted classifier (21
    features, 5-fold CV, AUC 0.853; the 14 largest shown).
    (c)~Distribution of \driftfeat{} over
    attacked runs, split by outcome: grey, runs without the harmful action;
    colored, runs with it (no Claude run had one); dashed line, the 0.75
    decision threshold.
    Per-violin run counts were not retained with this panel.}
  \label{fig:drift_appendix}
\end{figure}

\end{document}
